\documentclass[11pt]{article}

\usepackage[T1]{fontenc}
\usepackage[utf8]{inputenc}
\usepackage{lmodern}
\usepackage{geometry}
\usepackage{amsmath,amssymb,amsthm,mathtools}
\usepackage{booktabs}
\usepackage{microtype}
\usepackage{xurl}
\usepackage[hidelinks]{hyperref}
\geometry{margin=1in}

\newtheorem{theorem}{Theorem}[section]
\newtheorem{proposition}[theorem]{Proposition}
\newtheorem{corollary}[theorem]{Corollary}
\newtheorem{remark}[theorem]{Remark}
\newcommand{\R}{\mathbb{R}}
\newcommand{\rank}{\operatorname{rank}}
\newcommand{\Null}{\operatorname{ker}}
\newcommand{\diag}{\operatorname{diag}}

\title{A Structural Local Identifiability Obstruction\\
for the 31-Coefficient v18 Spectral Model}
\author{Richard J. Reyes}
\date{October 26, 2026}

\begin{document}
\maketitle

\begin{abstract}
We consider the eleven-nonzero-observable map of the v18 spectral
model with twenty-nine stored scalar coefficients and two historically
optimized parameters. At every differentiability point, rank--nullity
implies at least twenty first-order null directions. For the exact-decimal
frozen coefficients, $\epsilon=5$, and
$\alpha\in[5.172,5.173]$, we additionally prove that the actual clipped
31-parameter map is $C^1$ near the frozen vector, its derivative has rank
exactly eleven and nullity exactly twenty, and its local equal-output
fiber is a twenty-dimensional $C^1$ embedded submanifold. The proof uses
a selected minor of the full Fr\'echet derivative, a kernel-checked
trigonometric determinant enclosure, and the implicit-function theorem.
The theorem, proposition, and corollary statements below are formalized
in Lean 4 with Mathlib v4.23.0.
The historical finite-difference SVD remains a separate numerical
diagnostic. This is an audit of parameter identifiability, not an
independent physical prediction of particle masses.
\end{abstract}

\section{Model and hypotheses}

Let \(U\subseteq\R^{31}\) be an open neighborhood of a reference
coefficient vector \(\theta_0\), and let
\begin{equation}
  F:U\longrightarrow\R^{11},\qquad
  F(\theta)=\bigl(f_1(\theta),\ldots,f_{11}(\theta)\bigr)^{\mathsf T}
  \label{eq:forward}
\end{equation}
be the spectral forward map for the eleven nonzero observables.
The assumption that \(F\) is differentiable at \(\theta_0\) is explicit:
the implementation includes piecewise routing, clipping, and floors,
so differentiability should not be asserted at a branch transition.
Within a fixed smooth branch, the calculation below applies directly.
The rank--nullity conclusion also applies to any chosen linear
sensitivity matrix of the indicated shape, independently of how its
entries were estimated.

The audited parameter allocation is
\begin{center}
\begin{tabular}{lr}
\toprule
Parameter section & Number of scalar parameters\\
\midrule
Light sector & 9\\
Middle sector & 7\\
Heavy sector & 13\\
Historically live suppression parameters & 2\\
\midrule
Total & \(31\)\\
\bottomrule
\end{tabular}
\end{center}
The v18 source specifies twelve targets, including one artificial
zero-mass target. The full-coefficient audit excludes that zero target and
uses the index set
\begin{equation}
  \mathcal N=\{1,2,3,12,19,53,73,91,95,96,123\},
  \qquad |\mathcal N|=11.
  \label{eq:indices}
\end{equation}
The corresponding eleven outputs are evaluated without refitting the
stored coefficients or reassigning particle labels.

For the concrete certificate, $\theta_0$ is the vector
\texttt{V18Forward.frozen alpha 5}. Its twenty-nine stored coefficients
are the published finite decimals, interpreted as exact real numbers;
$\alpha\in[5.172,5.173]$. The Lean function \texttt{V18Forward.observed}
retains all thirty-one input coordinates, including both middle-sector
frequencies, both suppression parameters, all heavy powers, and the W/Z
frequency and phase. It includes the original floors, exponential and
power clipping, and routing. No equivalence with arbitrary NumPy
floating-point executions is claimed.


\section{Exact structural theorem}

\begin{theorem}[Minimum first-order nullity]
\label{thm:structural}
Suppose the map \(F\) in \eqref{eq:forward} is differentiable at
\(\theta_0\in U\). Its sensitivity matrix
\begin{equation}
  J=DF(\theta_0)
   =\left(\frac{\partial f_i}{\partial\theta_j}(\theta_0)\right)_
          {\substack{1\leq i\leq11\\1\leq j\leq31}}
   \in\R^{11\times31}
  \label{eq:jacobian}
\end{equation}
satisfies
\begin{equation}
  \boxed{\rank J\leq11},
  \qquad
  \boxed{\dim\Null J=31-\rank J\geq20}.
  \label{eq:main}
\end{equation}
In particular, there are at least twenty linearly independent directions
in parameter space along which all eleven predictions have vanishing
first directional derivative.
\end{theorem}

\begin{proof}
The derivative \(DF(\theta_0)\) is a linear map
\begin{equation}
 J:\R^{31}\longrightarrow\R^{11}.
\end{equation}
Because the dimension of its codomain is eleven,
\begin{equation}
 \rank J\leq\min\{31,11\}=11.
\end{equation}
The rank--nullity theorem gives
\begin{equation}
  \dim\Null J+\rank J=\dim\R^{31}=31,
\end{equation}
and hence
\begin{equation}
  \dim\Null J=31-\rank J\geq31-11=20.
\end{equation}
Choose linearly independent vectors
\(v_1,\ldots,v_d\) forming a basis of \(\Null J\), where \(d\geq20\).
For any \(v\in\Null J\), we have \(Jv=0\). Differentiability entails
\begin{equation}
 F(\theta_0+tv)
   =F(\theta_0)+tJv+o(|t|)
   =F(\theta_0)+o(|t|)
 \quad (t\to0).
 \label{eq:firstorder}
\end{equation}
Thus the first derivative of every component of \(F\) vanishes along
each \(v_i\), proving the last assertion.
\end{proof}

\begin{remark}[Meaning of ``unidentifiable direction'']
Equation \eqref{eq:firstorder} is a \emph{first-order} statement; it
neither says that a particular parameter is independently dispensable nor
that \(F(\theta_0+tv)=F(\theta_0)\) holds for nonzero \(t\).
The null vectors are generally combinations of the stored and live
coefficients. Furthermore, if the physically admissible parameter set
has active bounds, not every signed null perturbation need remain
feasible. The theorem concerns the ambient linearized map.
\end{remark}

\begin{corollary}[Exact nullity under exact full row rank]
\label{cor:exact}
If, in addition, \(\rank DF(\theta_0)=11\) holds \emph{exactly}, then
\begin{equation}
  \boxed{\dim\Null DF(\theta_0)=20}.
\end{equation}
\end{corollary}
\begin{proof}
Insert \(\rank J=11\) into rank--nullity.
\end{proof}

\section{Reproducing the numerical sensitivity audit}

The full-coefficient script loads the unchanged v18 implementation,
checks the original predictions against an independent explicit
branch-dispatch evaluator, and applies finite differences about the
historical parameter vector \(\theta_0\).
For a coordinate vector \(e_j\), it usually estimates
\begin{equation}
  J_{ij}\approx
   \frac{f_i(\theta_0+h_je_j)-f_i(\theta_0-h_je_j)}{2h_j},
  \qquad
  h_j=10^{-5}\max\{1,|\theta_{0,j}|\}.
  \label{eq:central}
\end{equation}
The historically live parameter \(\epsilon_{\mathrm{cutoff}}\)
is at its upper bound \(5\), so that column is computed using a
backward difference. Other active bounds, when relevant, also trigger
one-sided differences. These finite differences are diagnostics and
are not an exact derivative certificate.

Let \(y_i>0\) denote the eleven observed nonzero target masses and set
\(s_j=\max\{1,|\theta_{0,j}|\}\). The reported scaled matrix is
\begin{equation}
  \widetilde J= D_y^{-1} J D_s,
  \qquad
  D_y=\diag(y_1,\ldots,y_{11}),\quad
  D_s=\diag(s_1,\ldots,s_{31}).
  \label{eq:scaled}
\end{equation}
Both diagonal matrices are invertible, so exact linear-algebraic rank
is unchanged:
\begin{equation}
   \rank\widetilde J=\rank J.
\end{equation}
However, a numerical rank estimate can depend on scaling, finite
differences, and the threshold used in the SVD.

Writing the eleven singular values as
\(\sigma_1\geq\cdots\geq\sigma_{11}\geq0\), the script defines
\begin{equation}
  r_{\mathrm{num}}
   =\#\{i:\sigma_i>10^{-8}\sigma_1\}.
  \label{eq:rankthreshold}
\end{equation}
The executable audit reports
\begin{equation}
   \boxed{r_{\mathrm{num}}=11},\qquad
   \boxed{31-r_{\mathrm{num}}=20}.
\end{equation}
This is \emph{numerical} full row rank. By contrast, the lower bound
\(\dim\Null J\geq20\) in Theorem~\ref{thm:structural} is an exact
consequence of dimensions, even if the true rank were below eleven.

\section{Exact selected-minor certificate}
\label{sec:minor}

The proof selects the following zero-based coordinates of the complete
31-dimensional parameter vector, ordered by light, middle, and heavy sector:
\begin{equation}
 (0,3,4;\ 14,15,9,12,13;\ 20,24,28).
 \label{eq:selected}
\end{equation}
These are, respectively, $(m_0,\theta,\gamma)$, middle
$(c,d,A,\phi,h)$, and heavy (first nonlinear boson scale, W/Z-break
scale, Higgs offset). This selection differs from the earlier
finite-difference diagnostic's columns. It changes neither the forward
map nor its active parameters; it selects directions only \emph{after}
differentiating the complete map.

Strict inequalities at the frozen vector fix every relevant clipping
branch on an open neighborhood. All three light outputs and all five
middle outputs exceed their floors strictly. At the three heavy indices,
the inner light prediction is strictly below its floor, so that inner
term is locally the constant $10^{-12}$; each heavy outer prediction
is strictly above its floor. The relevant exponential and power clips
are also strictly inside, or strictly on the verified saturated side of,
their thresholds. Continuity propagates these inequalities locally.

Consequently, the selected restriction of the actual Fr\'echet derivative is
\begin{equation}
 J_{\mathrm{minor}}=
 \begin{pmatrix}L&0&0\\0&M&0\\0&0&H\end{pmatrix},
 \qquad \det J_{\mathrm{minor}}=\det L\,\det M\,\det H.
 \label{eq:block}
\end{equation}
All 121 entries, including all six off-diagonal blocks, are identified
in Lean. In particular, the heavy outputs' dependence on light coefficients
vanishes because of the verified inner floor, not because those coefficients
have been held fixed in the full derivative.

\subsection{Light sector}
The selected light block is $L=\diag(D_1,D_2,D_3)L_0$, where
\begin{equation}
 L_0=\begin{pmatrix}1/2&1&1\\2/3&4&8\\3/4&9&27\end{pmatrix},
 \qquad D_n=\frac{\eta_n}{1+\exp(\alpha(n-5))},
 \qquad (\eta_1,\eta_2,\eta_3)=(1,1,0.0048).
\end{equation}
Exact rational expansion gives
\begin{equation}
 \det L_0=9,\qquad \det L=9D_1D_2D_3>0.
 \label{eq:lightdet}
\end{equation}

\subsection{Middle sector}
The unfloored formula is
\begin{equation}
 f(n)=c+dn^2+h\sin(tn)+A\sin(sn+\phi).
\end{equation}
For $n\in\{12,19,53,73,91\}$, the selected derivative rows are
\begin{equation}
 M_n=(1,n^2,\sin(sn+\phi),A\cos(sn+\phi),\sin(tn)),
 \label{eq:middle}
\end{equation}
with $(A,s,t,\phi)=(2241.1,5.0632,9.4062,1.6469)$ at the frozen
vector. The frequencies remain active in the complete derivative.
The kernel-checked result for this exact matrix is
\begin{equation}
 \boxed{18{,}000{,}000<\det M<24{,}000{,}000}.
 \label{eq:middledet}
\end{equation}
The proof uses rational bounds on $\pi$, exact period reduction,
Mathlib sine and cosine Taylor remainder estimates, and rational interval
propagation through an exact 60-term determinant expansion. External
Python certificates are regression checks and supply no Lean premises.

\subsection{Heavy sector}
Let $a_n=(n/100)^{-1.189}$. The selected heavy block is
\begin{equation}
 H=\begin{pmatrix}
 a_{95}&\sin(-115.395)&0\\
 a_{96}&\sin(-116.659)&0\\
 a_{123}&0&1
 \end{pmatrix}.
\end{equation}
Here the first column differentiates the first nonlinear boson scale,
the second the W/Z-break amplitude, and the third the Higgs offset.
The full derivative also retains the power, frequency, and phase derivatives.
Period reduction with certified rational bounds on $\pi$ proves
$\sin(-115.395)<0<\sin(-116.659)$. Since $a_{95},a_{96}>0$,
\begin{equation}
 \det H=a_{95}\sin(-116.659)-a_{96}\sin(-115.395)>0.
 \label{eq:heavydet}
\end{equation}
Thus the actual selected derivative restriction is nonsingular. Its rank
is eleven and cannot exceed the rank of the full derivative. The latter
has eleven output coordinates, so
\begin{equation}
 \boxed{\rank DF(\theta_0)=11},\qquad
 \boxed{\dim\ker DF(\theta_0)=20}.
\end{equation}

\section{Local level-set geometry}

\begin{proposition}[Conditional nonlinear local level set]
\label{prop:levelset}
Assume \(F\) is continuously differentiable on an open neighborhood of
\(\theta_0\), and assume the \emph{exact} rank condition
\(\rank DF(\theta_0)=11\). Then, in a sufficiently small neighborhood
of \(\theta_0\), the level set
\begin{equation}
  \mathcal M=\{\theta:F(\theta)=F(\theta_0)\}
\end{equation}
is a \(C^1\) submanifold of dimension twenty.
\end{proposition}
\begin{proof}
Exact row rank eleven means that \(DF(\theta_0)\) is surjective.
The submersion (or implicit-function) theorem then provides local
coordinates in which \(F\) is the projection from \(\R^{31}\)
onto the first eleven coordinates, up to smooth changes of coordinates.
The fiber over \(F(\theta_0)\) is locally described by fixing those
eleven coordinates, leaving \(31-11=20\) free coordinates.
\end{proof}

\begin{corollary}[The concrete frozen fiber]
For the exact frozen vector with $\alpha\in[5.172,5.173]$ and
$\epsilon=5$, all hypotheses of Proposition~\ref{prop:levelset} hold.
More precisely, writing $K=\ker DF(\theta_0)$, there is an ambient
partial homeomorphism $\Psi$ between open neighborhoods in
$\R^{31}$ and $\R^{11}\times K$ such that
\begin{equation}
 \Psi(\theta_0)=(F(\theta_0),0),\qquad
 \pi_1\Psi(\theta)=F(\theta),\qquad \dim K=20.
\end{equation}
Both $\Psi$ and $\Psi^{-1}$ are $C^1$ on their open domains. The chart
maps the entire local fiber exactly onto
\[
 \{(y,k)\in\operatorname{target}(\Psi):y=F(\theta_0)\}.
\]
Every ambient neighborhood of $\theta_0$ contains a distinct parameter
vector with exactly the same eleven predictions.
\end{corollary}
\begin{proof}
The verified strict branches give $C^1$ regularity on an open
neighborhood. The certified minor proves surjectivity of the full
Fr\'echet derivative. Mathlib's implicit-function chart is restricted to
open sets where the chart and its inverse are $C^1$; the inverse
regularity follows from its inverse-function theorem. The first chart
coordinate is the actual observation map, and the second takes values
in its twenty-dimensional derivative kernel. Varying a nonzero kernel
coordinate slightly proves the final assertion. These steps are
exported by \texttt{V18Forward.f09\_completion} and
\texttt{V18Forward.observed\_locally\_nonidentifiable}.
\end{proof}


\begin{remark}[Active bounds]
The historical live parameter \(\epsilon_{\mathrm{cutoff}}\) sits at an
optimizer upper bound. Proposition~\ref{prop:levelset} is therefore
stated for a smooth extension to an \emph{open} ambient domain; it does
not automatically imply that every direction in the full twenty-dimensional
fiber respects the fitted parameter constraints. Since the certified
minor in Section~\ref{sec:minor} uses eleven
\emph{stored} parameters, it also permits a
restricted implicit-function analysis with the two live parameters fixed.
In that case the remaining twenty-nine stored parameters would support
a local \(29-11=18\)-dimensional equal-prediction level set, without
moving either live parameter across its bound.
\end{remark}

\section{Scope and scientific conclusion}

The theorem makes no claim that exactly twenty individually named
coefficients are redundant, that the measured masses identify the
remaining eleven coefficients, or that a fitted spectral formula follows
from a specified eigenvalue operator. It does not determine measurement
uncertainties or Bayesian posterior widths. In particular, it does not
supply an independent mapping from eigenmodes to particle assignments,
a confinement operator with its domain and boundary conditions, or a
physical normalization converting eigenvalues into measured mass units.

The logical status of the main statements is:
\begin{center}
\begin{tabular}{p{0.63\linewidth}p{0.29\linewidth}}
\toprule
Statement & Status\\
\midrule
\(\rank J\leq11\) for an 11-by-31 Jacobian & Exact theorem\\
\(\dim\Null J\geq20\) & Exact theorem\\
At least twenty zero first-order sensitivities & Exact when differentiable\\
The frozen audit's reported numerical rank eleven & Numerical diagnostic\\
Frozen rank eleven and nullity twenty & Kernel-checked exact result\\
Exactly twenty-dimensional smooth local fibers & Kernel-checked at the frozen vectors\\
Independent spectral prediction of particle masses & Not established\\
\bottomrule
\end{tabular}
\end{center}

\noindent\textbf{Conclusion.}
For the unrestricted 31-coefficient parameterization and eleven
nonzero observables, the first-order structural obstruction is
\begin{equation}
 \boxed{
  J\in\R^{11\times31}
  \quad\Longrightarrow\quad
  \rank J\leq11
  \quad\Longrightarrow\quad
  \dim\Null J\geq20.
 }
\end{equation}
The concrete frozen result additionally establishes a twenty-dimensional
nonlinear local fiber. It does not classify distant solutions, determine
the global topology of a fiber, or make a physically constrained parameter
box into an open manifold. None of these mathematical results constitutes
an independently validated physical spectrum. The corresponding v18 research
classification remains \texttt{EMPIRICAL\_FIT\_ONLY}.

\section*{Sources and reproducibility}
\begin{enumerate}
\item v18 audit and scientific-status discussion:\\
\url{https://github.com/rickyjreyes/geometry_of_resonance/blob/main/docs/SM_V18_SPECTRAL_BRIDGE_AUDIT.md}
\item Full-coefficient finite-difference/SVD audit:\\
\url{https://github.com/rickyjreyes/geometry_of_resonance/blob/main/SM_Spectrum/audit_v18_identifiability.py}
\item Original v18 source defining the coefficient blocks and predictions:\\
\url{https://github.com/rickyjreyes/geometry_of_resonance/blob/main/SM_Spectrum/_terms_perfection_v18.py}
\item Foundational Stage 2 spectral-bridge obstruction report:\\
\url{https://github.com/rickyjreyes/geometry_of_resonance/blob/main/docs/FOUNDATIONAL_STAGE2_2026-10-14.md}
\end{enumerate}

\noindent The repository's reproduction command is
\begin{verbatim}
python SM_Spectrum/audit_v18_identifiability.py \
  --output audit/v18_identifiability.json
\end{verbatim}
It records coefficient and CSV file hashes, package versions, the
11-by-31 scaled sensitivity, its SVD, and the declared numerical-rank
threshold. The historical numerical result is reported separately from the Lean
certificate above. No new executed holdout test is claimed.

The main formalization modules in \texttt{lean/MathProofs/} are:
\begin{itemize}
\item \texttt{V18MiddleDeterminant}: the exact trigonometric enclosure.
\item \texttt{V18Regularity} and \texttt{V18Frechet}: local branches and the full derivative.
\item \texttt{V18SelectedJacobian} and \texttt{V18Rank}: the actual minor, rank and nullity.
\item \texttt{LocalLevelSet} and \texttt{V18Identifiability}: the $C^1$ chart and completion theorem.
\end{itemize}
The aggregate command is
\begin{verbatim}
lake build MathProofs.V18Identifiability
\end{verbatim}
The complete repository is checked with \texttt{lake build}.

\end{document}
